\documentclass[aspectratio=169]{beamer}

% Blue-violet Beamer style inspired by the supplied reference
\usetheme{Madrid}
\usecolortheme{dolphin}
\definecolor{ConferenceBlue}{RGB}{52,54,183}
\definecolor{ConferenceBlueDark}{RGB}{42,44,155}
\setbeamercolor{background canvas}{bg=white}
\setbeamercolor{structure}{fg=ConferenceBlue}
\setbeamercolor{title}{fg=white,bg=ConferenceBlue}
\setbeamercolor{subtitle}{fg=white}
\setbeamercolor{frametitle}{fg=white,bg=ConferenceBlue}
\setbeamercolor{author in head/foot}{fg=white,bg=ConferenceBlue}
\setbeamercolor{title in head/foot}{fg=white,bg=ConferenceBlueDark}
\setbeamercolor{date in head/foot}{fg=white,bg=ConferenceBlue}
\setbeamercolor{page number in head/foot}{fg=white,bg=ConferenceBlueDark}
\setbeamercolor{section in toc}{fg=ConferenceBlue}
\setbeamerfont{title}{size=\LARGE}
\setbeamerfont{subtitle}{size=\large}
\setbeamerfont{author in head/foot}{size=\scriptsize}
\setbeamerfont{title in head/foot}{size=\scriptsize}
\setbeamerfont{date in head/foot}{size=\scriptsize}
\setbeamerfont{page number in head/foot}{size=\scriptsize}
\setbeamertemplate{footline}{%
  \leavevmode%
  \hbox{%
    \begin{beamercolorbox}[wd=.27\paperwidth,ht=2.8ex,dp=1.1ex,leftskip=.8em]{author in head/foot}%
      \usebeamerfont{author in head/foot}\insertshortauthor
    \end{beamercolorbox}%
    \begin{beamercolorbox}[wd=.36\paperwidth,ht=2.8ex,dp=1.1ex,center]{title in head/foot}%
      \usebeamerfont{title in head/foot}\insertshorttitle
    \end{beamercolorbox}%
    \begin{beamercolorbox}[wd=.25\paperwidth,ht=2.8ex,dp=1.1ex,center]{date in head/foot}%
      \usebeamerfont{date in head/foot}\insertshortdate
    \end{beamercolorbox}%
    \begin{beamercolorbox}[wd=.12\paperwidth,ht=2.8ex,dp=1.1ex,center]{page number in head/foot}%
      \usebeamerfont{page number in head/foot}\insertframenumber{} / \inserttotalframenumber
    \end{beamercolorbox}%
  }%
}
\setbeamertemplate{caption}[numbered]
\setbeamertemplate{bibliography item}[text]

\usepackage{booktabs}
\usepackage{tabularx}
\usepackage{array}
\usepackage{amsmath}

% Replace each gray placeholder with the corresponding paper content
\newcommand{\fillin}[1]{\textcolor{gray}{\textit{[#1]}}}

% Screenshot-style title page with a rounded blue title banner
\setbeamertemplate{title page}{%
  \begin{centering}
    \vspace*{0.35cm}
    \begin{beamercolorbox}[wd=.96\paperwidth,sep=1.2em,center,rounded=true,shadow=true]{title}
      {\usebeamerfont{title}\inserttitle\par}
      \vspace{0.3em}
      {\usebeamerfont{subtitle}\insertsubtitle\par}
    \end{beamercolorbox}
    \vspace{1.0cm}
    {\usebeamerfont{author}\insertauthor\par}
    \vspace{0.55cm}
    {\usebeamerfont{institute}\insertinstitute\par}
    \vfill
    {\usebeamerfont{date}\insertdate\par}
  \end{centering}
}

% ===== Presentation details =====
\title[Research Talk]{[Academic Presentation Title]}
\subtitle{[Presentation Subtitle]}
\author[Presenter, Co-author]{[Presenter Name]\inst{1}\hspace{1em}[Co-author Name]\inst{2}}
\institute[Institution]{%
  \inst{1}[Department / Institution]\\
  \inst{2}[Department / Institution]
}
\date[Conference, Month Year]{[Conference Name, Month Year]}
\begin{document}

% Slide 1: Title
\begin{frame}
  \titlepage
\end{frame}

% Slide 2: Definition of progressive visualization
\begin{frame}{Progressive vs. Non-progressive Visualization}
  \begin{columns}[T,onlytextwidth]
    \begin{column}{0.46\textwidth}
      \begin{block}{Non-progressive (``eager'') visualization}
        \small
        \begin{itemize}
          \setlength{\itemsep}{0.65em}
          \item Renders in one pass.
          \item A data change or user interaction triggers a full rendering.
        \end{itemize}
      \end{block}
    \end{column}
    \begin{column}{0.52\textwidth}
      \begin{block}{Progressive visualization}
        \small
        \begin{itemize}
          \setlength{\itemsep}{0.55em}
          \item Renders in a sequence of passes, progressively improving toward completion.
          \item Each pass stays within a requested latency and produces a meaningful partial result.
          \item Each partial result includes an error or quality measure and an estimate of completion.
          \item Allows interaction at any time, adapts when possible, and can support steering.
        \end{itemize}
      \end{block}
    \end{column}
  \end{columns}
  \vfill
  \centering
  \small\textbf{Meaningful partial results are part of the definition.}\\[0.35em]
  {\scriptsize\textcolor{gray}{Definition adapted from Ulmer et al. (2024), Sec.~2, ``Formal Definition,'' PDF p.~2.}}
\end{frame}

\section{Literature Review}

% Review slide 1 of 6
\begin{frame}{Motivation and Scope}
  \begin{itemize}
    \item Larger datasets and longer-running analytics delay the first usable result and interrupt visual exploration.
    \item Progressive visualization shows meaningful partial views while computation continues; the views improve over time.
    \item Ulmer et al. focus on visual and interactive systems, excluding work that only accelerates algorithms without progressive visual output.
  \end{itemize}
  \vfill
  {\scriptsize\textcolor{gray}{Source: Ulmer et al. (2024).}}
\end{frame}

% Review slide 2 of 6
\begin{frame}{Requirements for Progressive Visualization}
  \begin{itemize}
    \item Computation proceeds in bounded passes, each returning a meaningful partial result.
    \item Partial results should communicate quality or error and give an estimate of completion.
    \item Users can interact before completion; steering can direct further computation toward relevant regions.
    \item Data chunking addresses large inputs; process chunking exposes progress in slow algorithms.
  \end{itemize}
  \vfill
  {\scriptsize\textcolor{gray}{Sources: Ulmer et al. (2024); Pezzotti et al. (2017).}}
\end{frame}

% Review slide 3 of 6
\begin{frame}{A Taxonomy of Progressive Visualizations}
  \begin{itemize}
    \item \textbf{Visualization type:} temporal, geo-spatial, hierarchical, network, multidimensional, or field.
    \item \textbf{Processing:} data chunking, process chunking, or custom chunking.
    \item \textbf{Data domain:} known end or unknown end.
    \item \textbf{Visual update:} extend the current view or overwrite it with a new state.
  \end{itemize}
  \vfill
  {\scriptsize\textcolor{gray}{Source: Ulmer et al. (2024).}}
\end{frame}

% Review slide 4 of 6
\begin{frame}{Progressive Dimensionality Reduction}
  \begin{itemize}
    \item t-SNE builds an embedding through iterative optimization, but it must first compute high-dimensional similarities.
    \item Barnes-Hut t-SNE reduces optimization cost, yet initialization can still delay the first interpretable embedding.
    \item Q-SNE approximated neighbors for sparse document data; A-tSNE extends approximate similarities to dense high-dimensional data.
  \end{itemize}
  \vfill
  {\scriptsize\textcolor{gray}{Source: Pezzotti et al. (2017), Related Work.}}
\end{frame}

% Review slide 5 of 6
\begin{frame}{Field Trends and Open Challenges}
  \begin{itemize}
    \item The survey classifies 48 papers from 1999 to 2023; 33 of 48 use multidimensional visualization.
    \item Steering appears in 35 studies, uncertainty in 20, and real-time processing in 10.
    \item Open challenges include visual stability as estimates change, uncertainty communication, and evaluation beyond runtime.
  \end{itemize}
  \vfill
  {\scriptsize\textcolor{gray}{Source: Ulmer et al. (2024), study set and discussion.}}
\end{frame}

% Review slide 6 of 6
\begin{frame}{Synthesis and Paper Positioning}
  \begin{itemize}
    \item A-tSNE addresses an algorithm-level bottleneck: delayed initialization in high-dimensional t-SNE.
    \item The 2024 survey offers a field-level taxonomy and reviews evaluation methods and research gaps.
    \item \textbf{Synthesis:} progressive methods must make partial results timely, interpretable, steerable, and stable as they improve.
  \end{itemize}
  \vfill
  {\scriptsize\textcolor{gray}{Sources: Pezzotti et al. (2017); Ulmer et al. (2024).}}
\end{frame}
\section{Paper 1}

% Slides 9-18: Paper 1 (10 slides)
\begin{frame}{Paper 1: Overview}
  \begin{itemize}
    \item \fillin{Title, authors, and publication details}
    \item \fillin{Research subject and context}
    \item \fillin{Core question addressed in this paper}
  \end{itemize}
\end{frame}

\begin{frame}{Paper 1: Motivation}
  \begin{itemize}
    \item \fillin{Why the research question matters}
    \item \fillin{Limitations of existing explanations}
    \item \fillin{How the paper approaches the question}
  \end{itemize}
\end{frame}

\begin{frame}{Paper 1: Research Questions and Hypotheses}
  \begin{itemize}
    \item \fillin{Research question}
    \item \fillin{Hypothesis or proposition 1}
    \item \fillin{Hypothesis or proposition 2, if applicable}
  \end{itemize}
\end{frame}

\begin{frame}{Paper 1: Conceptual Framework}
  \begin{itemize}
    \item \fillin{Key concepts or variables}
    \item \fillin{Theoretical relationships and proposed mechanisms}
    \item \fillin{Insert the paper's central framework figure here}
  \end{itemize}
\end{frame}

\begin{frame}{Paper 1: Data and Research Subjects}
  \begin{itemize}
    \item \fillin{Data or material sources and time period}
    \item \fillin{Sample composition or case selection}
    \item \fillin{Data processing and ethics, if applicable}
  \end{itemize}
\end{frame}

\begin{frame}{Paper 1: Research Methods}
  \begin{itemize}
    \item \fillin{Research design and analytical method}
    \item \fillin{Model, experiment, or analysis workflow}
    \item \fillin{Key identification assumptions or derivation conditions}
  \end{itemize}
\end{frame}

\begin{frame}{Paper 1: Results Overview}
  \begin{itemize}
    \item \fillin{Main result 1}
    \item \fillin{Main result 2}
    \item \fillin{How the results address the hypotheses or propositions}
  \end{itemize}
\end{frame}

\begin{frame}{Paper 1: Main Evidence}
  \begin{itemize}
    \item \fillin{Insert the key figure, table, or theorem}
    \item \fillin{Key estimate, proof, or observation}
    \item \fillin{Conclusion supported by this evidence}
  \end{itemize}
\end{frame}

\begin{frame}{Paper 1: Mechanisms and Additional Analyses}
  \begin{itemize}
    \item \fillin{Mechanism analysis or theoretical explanation}
    \item \fillin{Robustness, sensitivity, or alternative explanation checks}
    \item \fillin{Boundary conditions for the results}
  \end{itemize}
\end{frame}

\begin{frame}{Paper 1: Summary}
  \begin{itemize}
    \item \fillin{Main conclusion}
    \item \fillin{Contribution and limitation}
    \item \fillin{Connection to Paper 2}
  \end{itemize}
\end{frame}

\section{Paper 2}

% Slides 19-26: Paper 2 (8 slides)
\begin{frame}{Paper 2: Overview}
  \begin{itemize}
    \item \fillin{Title, authors, and publication details}
    \item \fillin{Research subject and context}
    \item \fillin{Core question addressed in this paper}
  \end{itemize}
\end{frame}

\begin{frame}{Paper 2: Research Question and Contribution}
  \begin{itemize}
    \item \fillin{Research question}
    \item \fillin{What this paper adds to related work}
    \item \fillin{Relationship to or distinction from Paper 1}
  \end{itemize}
\end{frame}

\begin{frame}{Paper 2: Framework and Hypotheses}
  \begin{itemize}
    \item \fillin{Theoretical basis or conceptual framework}
    \item \fillin{Hypotheses or propositions}
    \item \fillin{Proposed mechanism or reasoning}
  \end{itemize}
\end{frame}

\begin{frame}{Paper 2: Data and Research Design}
  \begin{itemize}
    \item \fillin{Data or material sources and sample}
    \item \fillin{Research design and key variables or concepts}
    \item \fillin{Why the design fits the research question}
  \end{itemize}
\end{frame}

\begin{frame}{Paper 2: Methods and Analysis Strategy}
  \begin{itemize}
    \item \fillin{Model, experiment, comparison, or proof method}
    \item \fillin{Identification strategy or analytical steps}
    \item \fillin{Key assumptions and how they are assessed}
  \end{itemize}
\end{frame}

\begin{frame}{Paper 2: Main Results}
  \begin{itemize}
    \item \fillin{Main result 1 and supporting evidence}
    \item \fillin{Main result 2 and supporting evidence}
    \item \fillin{How the results answer the research question}
  \end{itemize}
\end{frame}

\begin{frame}{Paper 2: Mechanisms and Robustness}
  \begin{itemize}
    \item \fillin{Mechanisms, heterogeneity, or boundary conditions}
    \item \fillin{Robustness or sensitivity analyses}
    \item \fillin{Alternative explanations and additional evidence}
  \end{itemize}
\end{frame}

\begin{frame}{Paper 2: Summary and Overall Conclusions}
  \begin{itemize}
    \item \fillin{Main conclusion of Paper 2}
    \item \fillin{Shared findings and differences across the two papers}
    \item \fillin{Overall answer to the research topic}
  \end{itemize}
\end{frame}

% Slide 27: References
\section{References}
\begin{frame}{References}
  \small
  \begin{thebibliography}{99}
    \bibitem{pezzotti2017} N. Pezzotti et al., Approximated and User Steerable t-SNE for Progressive Visual Analytics, \textit{IEEE Transactions on Visualization and Computer Graphics}, 23(7), 1739--1752, 2017. doi:10.1109/TVCG.2016.2570755.

    \bibitem{ulmer2024} A. Ulmer et al., A Survey on Progressive Visualization, \textit{IEEE Transactions on Visualization and Computer Graphics}, 30(9), 6447--6467, 2024. doi:10.1109/TVCG.2023.3346641.
\end{thebibliography}
\end{frame}

% Slide 28: Questions
\begin{frame}
  \centering
  \vfill
  {\LARGE \fillin{Thank You}}\\[1em]
  {\large \fillin{Questions and Discussion}}\\[2em]
  \fillin{Presenter email, optional}
  \vfill
\end{frame}

\end{document}







